\documentclass[a4paper,12pt]{article}
\usepackage[slovene]{babel}
\usepackage[utf8]{inputenc}
\usepackage[T1]{fontenc}
\usepackage{lmodern}
\usepackage{amsmath,amsfonts}
%\usepackage{enumitem}
%\usepackage{mathabx}
\usepackage{fancyhdr}
\usepackage{url}
\usepackage{graphicx}
\usepackage{xcolor}
\usepackage{amsthm}
\usepackage{amssymb}

\theoremstyle{definition}
\newtheorem{definicija}{Definicija}

\title{Metric dimension of $Cn(1, 2, . . . , n - d)$: exact values}
\author{Živa Hegler, Neža Zupančič, Sara Polajnar}

\begin{document}
\maketitle

\section{Osnovne definicje}
V vseh nadaljnih definicijah predpostavimo, da imamo usmerjen graf
$G = (V, E)$, kjer je $V$ množica vseh vozlišč in $E$ množica vseh usmerjenih povezav. 
\begin{definicija}
    Naj bodo $1 \leq a_1 < a_2 \ldots < a_k \leq n-1$ generatorji.
    Usmerjeni cirkulantni graf $C_n(a_1, \ldots, a_{k})$ ima vozlišča
    $v_0, v_1, \ldots, v_{n-1}$ in usmerjene povezave $v_iv_{i+a_j}$ za  
    $0 \leq i \leq n-1$ in $1 \leq j \leq k$, pri čemer indekse računamo
    po modulu n. V našem projektu smo obravnavale usmerjene cirkulantne grafe 
    $C_n(1,2,\ldots,n-d)$, kjer je $n$ število vozlišč in $n-d$ število generatorjev.
\end{definicija}

\begin{definicija}
    Naj bo $W = \{w_1, \ldots, w_z\} \subseteq V$ množica vozlišč.
    Za vsako vozlišče $v \in V$ definiramo vektor razdalj glede na
    množico $W$ kot:
    $$
    r(W|v) = (d(w_1, v), \ldots, d(w_z, v)),
    $$
    kjer $d(w_i, v)$ označuje najkrajšo dolžino usmerjen poti med
    vozliščem $w_i$ in vozliščem $v$. Množiva $W$ je \textbf{razločevalna množica},
    če za vsaki različni vozlišči $u, v \in V$ velja:
    $$
    r(W|u) \neq r(W|v)
    $$
\end{definicija}

\begin{definicija}
    \textbf{Metrična dimenzija} grafa $G$ je moč najmanjše razločevalne
    množice. Označimo jo z $dim(G)$
\end{definicija}


\section{Načrt dela}
\subsection{Naloga 1: Implementacija ILP modela in testiranje na znanih vrednostih}

V prvem koraku bomo implementirale program za natančen izračun metrične dimenzije
$\operatorname{dim}(C_n(1, \dots, n-d))$ in iskanje pripadajoče optimalne razrešujoče množice $W$. 
Za osnovo bomo vzele ILP model, ki sta ga v svoji projektni nalogi uporabili že Jasmina Mazej in Lana Stojčić (avtorica 
kode za omenjen ILP je Maša Popovič), ter ga ustrezno prilagodile specifikam našega problema.

Ker se omejujemo na območje $n > 2d$, opazimo, da ima graf naslednjo lastnost: razdalja med poljubnima različnima 
vozliščema $v_i$ in $v_j$ je vedno bodisi 1 bodisi 2:
\[
d(v_i, v_j) = 
\begin{cases}
0, & \text{če } i = j, \\
1, & \text{če } (j - i) \bmod n \in \{1, 2, \dots, n-d\}, \\
2, & \text{sicer},
\end{cases}
\]
kjer z $(j - i) \bmod n \in \{0, 1, \dots, n-d\}$ označujemo najmanjši nenegativni ostanek pri deljenju z $n$, kar 
geometrijsko ustreza številu korakov od $v_i$ do $v_j$ v smeri urinega kazalca.

Problem formuliramo kot problem pokritja (angl.~\emph{Hitting Set}):
\begin{itemize}
    \item Za vsako vozlišče $w \in V$ uvedemo binarno spremenljivko $x_w \in \{0, 1\}$, kjer $x_w = 1$ pomeni, da je $w \in W$.
    \item Za vsak neurejen par različnih vozlišč $\{u, v\}$ določimo množico vozlišč, ki ta par ločijo:
    \[
    D_{u,v} = \{w \in V \mid d(w, u) \neq d(w, v)\}.
    \]
    \item Zahtevamo, da je vsak par $\{u, v\}$ ločen z vsaj enim izbranim vozliščem:
    \[
    \sum_{w \in D_{u,v}} x_w \geq 1 \quad \text{za vse } u < v.
    \]
    \item Ciljna funkcija minimizira velikost razrešujoče množice: $\min \sum_{w \in V} x_w$.
\end{itemize}

\paragraph{Preverjanje pravilnosti implementacije:}
Pravilnost programa bomo sistematično preverile s primerjavo izračunanih rezultatov z že dokazanimi formulami iz 
literature:

\begin{itemize}
    \item Za $d = 4$ in $n \geq 10$:
    \[
    \operatorname{dim}(C_n(1, \dots, n-4)) = \left\lceil \frac{2n}{5} \right\rceil.
    \]
    \item Za $d = 5$ in $n \geq 11$:
    \[
    \operatorname{dim}(C_n(1, \dots, n-5)) = 
    \begin{cases}
    \lfloor n/3 \rfloor, & \text{če } 3 \mid n, \\
    \lfloor n/3 \rfloor + 2, & \text{če } 3 \nmid n.
    \end{cases}
    \]
\end{itemize}

\subsection{Naloga 2: Računanje metrične dimenzije za večje grafe in optimizacija reševalnika}

Ko bo program preverjen na manjših primerih, bomo izvedle sistematične računske eksperimente za parametre v območju:
\[
4 \leq d \leq 15 \quad \text{in} \quad n > 2d.
\]
Cilj je izračunati točne vrednosti metrične dimenzije čim dlje prek meje $n = 40$, pri čemer bomo za parametre 
$5 \leq d \leq 10$ rezultate izračunale vsaj do $n = 100$.

\paragraph{Obvladovanje časovne zahtevnosti:}
Navodila opozarjajo, da privzeti reševalnik v SageMath pri nekaterih parametrih postane opazno počasnejši. Za reševanje 
tega problema bomo uporabile naslednji pristop:
\begin{itemize}
    \item \textbf{Fiksiranje začetnega vozlišča:} Cirkulantni grafi so invariantni za ciklične zamike vozlišč. 
    Če obstaja kakšna optimalna razrešujoča množica, jo lahko z rotacijo vedno preslikamo tako, da vsebuje vozlišče 
    $v_0$.
\end{itemize}

\paragraph{Zbiranje podatkov:}
Vse dobljene rezultate bomo sproti shranjevale v tabelarični obliki (datoteke CSV). Za vsak par $(n, d)$ bomo zabeležile:
\begin{itemize}
    \item izračunano metrično dimenzijo $\operatorname{dim}(C_n(1, \dots, n-d))$,
    \item najdeno optimalno razrešujočo množico $W$,
    \item čas izvajanja računanja (za analizo časovne zahtevnosti).
\end{itemize}

\subsection{Naloga 3: Testiranje Domneve 1 in pogojev za enakost}

Tretja faza projekta je namenjena empiričnemu preverjanju \textbf{Domneve 1}, ki sta jo v navodilih zastavila profesorja. 
Domneva podaja teoretično spodnjo mejo za metrično dimenzijo in napoveduje pogoje, kdaj je ta meja dosežena z enakostjo.

Za $d \geq 4$ in $n > 2d$ je predlagana spodnja meja definirana kot:
\[
\operatorname{dim}(C_n(1, \dots, n-d)) \geq B(n, d) = 
\begin{cases}
\left\lceil \frac{n}{d+1} \right\rceil \left\lceil \frac{d+1}{3} \right\rceil, & \text{če } d \not\equiv 0 \pmod 3, \\[1.5ex]
\frac{n}{2d-1} \cdot \frac{2d}{3}, & \text{če } d \equiv 0 \pmod 3.
\end{cases}
\]
Ker je metrična dimenzija celo število, se v primeru $d \equiv 0 \pmod 3$ desna stran zaokroži navzgor na celo število $\lceil B(n, d) \rceil$.

V sklopu te naloge bomo na vseh izračunanih podatkih  preverile:
\begin{itemize}
    \item \textbf{Veljavnost spodnje meje:} Preverile bomo, ali za vse pare $(n, d)$ v našem naboru podatkov velja $\operatorname{dim} \geq \lceil B(n, d) \rceil$, torej ali morda kje obstaja kakšen protiprimer.
    \item \textbf{Karakterizacija enakosti:} Domneva trdi, da je meja dosežena (enakost $\operatorname{dim} = \lceil B(n, d) \rceil$) natanko tedaj, ko sta izpolnjena naslednja pogoja deljivosti:
    \begin{itemize}
        \item za $d \not\equiv 0 \pmod 3$: enakost velja $\iff (d+1) \mid n$,
        \item za $d \equiv 0 \pmod 3$: enakost velja $\iff (2d-1) \mid n$.
    \end{itemize}
    Preverile bomo obe smeri ekvivalence: ali enakost res nastopi vedno, ko je pogoj deljivosti izpolnjen, in ali 
    vedno odpove, ko pogoj ni izpolnjen.
    \item \textbf{Analiza odstopanja:} Za primere, ko enakost ne velja, bomo spremljale razliko med dejansko dimenzijo 
    in spodnjo mejo ($\operatorname{dim} - \lceil B(n, d) \rceil$).
\end{itemize}

\subsection{Naloga 4: Iskanje eksaktnih formul za $d \geq 6$}

Zadnji cilj projekta je na podlagi zbranih empiričnih podatkov poskusiti uganiti zaprte eksaktne formule za metrično 
dimenzijo pri posameznih fiksnih vrednostih $d = 6, 7, 8, \dots$, podobno kot sta v literaturi že znani formuli za 
$d = 4$ in $d = 5$.

\paragraph{Postopek analize in iskanja vzorcev:}
Ker se vedenje grafov tesno navezuje na pogoje enakosti iz Domneve~1, bomo podatke za vsak fiksni $d$ razdelile v 
ostankarske razrede glede na ustrezen modul:
\begin{itemize}
    \item po modulu $d + 1$, kadar velja $d \not\equiv 0 \pmod 3$,
    \item po modulu $2d - 1$, kadar velja $d \equiv 0 \pmod 3$.
\end{itemize}
Če zapišemo $n = q \cdot M + r$ (kjer je $M$ ustrezen modul, $r \in \{0, 1, \dots, M-1\}$ ostanek, $q$ pa celi del 
količnika), bomo za vsak ostanek $r$ posebej opazovale, kako se dimenzija spreminja z naraščanjem $q$. Iskale bomo 
stabilne linearne vzorce oblike:
\[
\operatorname{dim} = a \cdot q + c_r
\]
oziroma ekvivalentne zapise s funkcijama stropa $\lceil \cdot \rceil$ in dna $\lfloor \cdot \rfloor$.

\paragraph{Določitev začetne meje veljavnosti:}
Pri cirkulantnih grafih se pri majhnih $n$ pogosto pojavijo začetna odstopanja, preden se vzorec ustali. Zato bomo za 
vsako predlagano formulo določile:
\begin{itemize}
    \item prag $N_d$ (najmanjši $n$), od katerega naprej predlagana formula velja za vse testirane primere brez izjeme,
    \item natančen seznam vseh odstopanj (anomalij) pod to mejo ($n < N_d$).
\end{itemize}

Vse tako dobljene formule bomo v končnem poročilu predstavile s tabelami in grafi ter jih jasno označile kot empirično 
podprte domneve.
\end{document}
